\subsubsection{Spectral packing}
\label{hb:spectral-subsection}

We first prove the geometric estimate underlying the construction of
bags.  It is the spectral packing statement of
\cite[Lemma~6.4]{AOG}; the argument below gives a logarithmic bound and
chooses the radius explicitly.
Here a Euclidean ball of \emph{squared radius} \(r\) means
\[
 \mathcal B_2(y,r)=\{z:\|z-y\|_2^2\le r\}.
\]
In particular, two such closed balls are disjoint when their centers
are at distance strictly greater than \(2\sqrt r\).

\begin{lemma}[Spectral packing]\label{hb:spectral-packing}
Let \(F\) be a nonempty set of oriented edges of a finite multigraph
with edge set \(E\), and put \(m=|F|\).
For points \(Y_v\in\R^E\), write
\(Y_e=Y_u-Y_v\) when \(e=(u,v)\), and denote the \(e\)-coordinate of
\(Y_e\) by \(Y_{e,e}\).
Suppose that \(a>0\) and
\begin{equation}\label{hb:spectral-assumption}
 \Upsilon:=\left(\frac1m\sum_{e\in F}Y_{e,e}\right)^2>0,
 \qquad
 \Upsilon\ge\frac a m\sum_{e\in F}\|Y_e\|_2^2.
\end{equation}
Then \(a\le1\), and there are \(b\ge2\) pairwise disjoint Euclidean
balls of one positive squared radius \(r\), centered at images of
endpoints of edges in \(F\), such that
\begin{equation}\label{hb:spectral-logarithmic}
 b\ge\frac{am}{16},
 \qquad
 br\ge\frac{9m\Upsilon}
 {512\bigl(5+\log_2(1/a)\bigr)^2}.
\end{equation}
Consequently, for every \(0<\varepsilon<1/3\), the choice
\begin{equation}\label{hb:spectral-constant}
 C_1(\varepsilon)=\frac{1024}{\varepsilon^2}
\end{equation}
gives
\begin{equation}\label{hb:spectral-conclusion}
 b\ge\frac{a|F|}{C_1(\varepsilon)},
 \qquad
 br\ge\frac{a^\varepsilon\Upsilon|F|}{C_1(\varepsilon)}.
\end{equation}
\end{lemma}

\begin{proof}
The scalar Cauchy--Schwarz inequality gives
\[
 \Upsilon\le\frac1m\sum_{e\in F}Y_{e,e}^2
 \le\frac1m\sum_{e\in F}\|Y_e\|_2^2.
\]
The last quantity is positive, so the assumption implies \(a\le1\).

\smallskip
\noindent\emph{The matrix estimate.}
Let \(P_F:\R^E\to\R^F\) be coordinate restriction, and form the
\(m\)-by-\(m\) matrix \(A\) whose column indexed by \(e\in F\)
is \(P_FY_e\).  Set \(\tau=m\sqrt\Upsilon\).
The singular values \(\sigma_1\ge\cdots\ge\sigma_m\ge0\) of \(A\)
satisfy
\begin{equation}\label{hb:spectral-matrix-bounds}
 \sum_{i=1}^m\sigma_i\ge|\tr A|=\tau,
 \qquad
 \sum_{i=1}^m\sigma_i^2=\|A\|_F^2
 \le\frac{\tau^2}{am}.
\end{equation}
Indeed, writing a singular value decomposition as
\(A=\sum_i\sigma_i u_i v_i^{\mathsf T}\), the trace is
\(\sum_i\sigma_i\langle u_i,v_i\rangle\), whose absolute value
is at most \(\sum_i\sigma_i\).  The Frobenius norm identity follows
from orthonormality, and its upper bound follows from
\eqref{hb:spectral-assumption} and the fact that coordinate restriction
does not increase a vector's norm.

We will also use the following rank approximation inequality:
for \(0\le k\le m\) and any matrix \(C\) of rank at most \(k\),
\begin{equation}\label{hb:spectral-rank-bound}
 \|A-C\|_F^2\ge\sum_{i=k+1}^m\sigma_i^2.
\end{equation}
For completeness, let \(Q\) be orthogonal projection onto the column
space of \(C\).  Since \((I-Q)C=0\), orthogonal projection and the
singular value decomposition give
\[
 \|A-C\|_F^2
 \ge\|(I-Q)A\|_F^2
 =\sum_{i=1}^m\sigma_i^2\bigl(1-\|Qu_i\|_2^2\bigr).
\]
The numbers \(w_i=\|Qu_i\|_2^2\) lie in \([0,1]\), and
\(\sum_iw_i=\rank Q\le k\).
Because the \(\sigma_i^2\) are nonincreasing,
\(\sum_i\sigma_i^2w_i\le\sum_{i=1}^k\sigma_i^2\):
moving available weight from a later index to an earlier index never
decreases this weighted sum, and its maximum therefore puts weight one
on the first \(k\) indices.  This proves
\eqref{hb:spectral-rank-bound}, including \(k=0\).

\smallskip
\noindent\emph{Selecting a singular-value tail.}
Define
\[
 p_0=\left\lfloor\frac{am}{16}\right\rfloor+1,
 \qquad T_p=\sum_{i=p}^m\sigma_i^2\quad(1\le p\le m).
\]
We have \(1\le p_0\le m\).  By \eqref{hb:spectral-matrix-bounds},
\[
 \sum_{i<p_0}\sigma_i
 \le\sqrt{(p_0-1)\sum_i\sigma_i^2}
 \le\frac\tau4.
\]
The sum is zero when \(p_0=1\), so this also covers arbitrarily small
\(am\).  It follows that
\begin{equation}\label{hb:spectral-surviving-sum}
 \sum_{i=p_0}^m\sigma_i\ge\frac{3\tau}{4}.
\end{equation}
Partition the integer interval \([p_0,m]\) into the nonempty blocks
\[
 I_j=\{2^jp_0,\ldots,\min(2^{j+1}p_0-1,m)\},
 \qquad 0\le j<L,
\]
where
\[
 L=1+\left\lfloor\log_2\frac m{p_0}\right\rfloor
 \le 5+\log_2(1/a).
\]
For the first index \(p=2^jp_0\) of a block, its cardinality is at
most \(p\).  Cauchy--Schwarz therefore gives
\[
 \sum_{i\in I_j}\sigma_i\le\sqrt{pT_p}.
\]
Summing this inequality and using
\eqref{hb:spectral-surviving-sum}, one of these first indices \(p\)
satisfies
\begin{equation}\label{hb:spectral-large-tail}
 p\ge p_0>\frac{am}{16},
 \qquad
 pT_p\ge\frac{9\tau^2}{16L^2}>0.
\end{equation}

\smallskip
\noindent\emph{Selecting the balls at an explicit radius.}
Fix the positive number
\begin{equation}\label{hb:spectral-radius}
 r=\frac{T_p}{32m}.
\end{equation}
Scan the endpoint images of edges in \(F\), keeping a point precisely
when its distance from every previously kept point is greater than
\(2\sqrt r\).  This finite procedure produces centers
\(Y_{w_1},\ldots,Y_{w_b}\) for pairwise disjoint closed balls of
squared radius \(r\).  Its maximality also ensures that each endpoint
image \(Y_v\) lies within distance \(2\sqrt r\) of one of the
centers; choose such a center and call its vertex \(c(v)\).

If \(b\le p\), define a matrix \(C\) by
\[
 C_e=P_F\bigl(Y_{c(u)}-Y_{c(v)}\bigr)
 \qquad(e=(u,v)\in F).
\]
Every column of \(C\) belongs to the span of
\(P_F(Y_{w_j}-Y_{w_1})\), \(2\le j\le b\), so
\(\rank C\le b-1\le p-1\).  On the other hand,
\begin{align*}
 \|A-C\|_F^2
 &\le\sum_{e=(u,v)\in F}
 \bigl\|(Y_u-Y_{c(u)})-(Y_v-Y_{c(v)})\bigr\|_2^2\\
 &\le\sum_{e=(u,v)\in F}
 2\bigl(\|Y_u-Y_{c(u)}\|_2^2+
          \|Y_v-Y_{c(v)}\|_2^2\bigr)
 \le16mr=\frac{T_p}{2}.
\end{align*}
This contradicts \eqref{hb:spectral-rank-bound} with \(k=p-1\),
which bounds the same quantity below by \(T_p>0\).
Thus \(b\ge p+1\ge2\).  Combining
\eqref{hb:spectral-large-tail} and \eqref{hb:spectral-radius} yields
\[
 b>\frac{am}{16},
 \qquad
 br\ge\frac{pT_p}{32m}
 \ge\frac{9m\Upsilon}{512L^2}.
\]
The bound on \(L\) proves \eqref{hb:spectral-logarithmic}.

\smallskip
\noindent\emph{Conversion to the power bound.}
Put \(t=\log(1/a)\ge0\).  Since \(\log2\ge1/2\),
\(5+\log_2(1/a)\le5+2t\).
The elementary maximum of \(t e^{-\varepsilon t/2}\) is
\(2/(e\varepsilon)\le1/\varepsilon\), and
\(5\le5/(3\varepsilon)\) for \(0<\varepsilon<1/3\).
Consequently,
\[
 \bigl(5+\log_2(1/a)\bigr)e^{-\varepsilon t/2}
 \le5+\frac2\varepsilon<\frac4\varepsilon.
\]
Squaring and substituting in \eqref{hb:spectral-logarithmic} gives
\[
 br\ge\frac{9\varepsilon^2}{8192}
       a^\varepsilon m\Upsilon
 \ge\frac{\varepsilon^2}{1024}
       a^\varepsilon m\Upsilon.
\]
Finally \(C_1(\varepsilon)\ge16\), so the count estimate in
\eqref{hb:spectral-conclusion} follows as well.
\end{proof}

The positivity assumption on \(\Upsilon\) is needed when one asks for
many distinct centers.  If \(\Upsilon=0\),
\eqref{hb:spectral-assumption} forces \(Y_e=0\) for every \(e\in F\),
and all endpoint images may coincide.  Such a set contributes zero to
the squared sums for which this lemma is used and can be discarded.
